%%%% ijcai26.tex -- ProtBLIP2-SST

\typeout{IJCAI--ECAI 26 ProtBLIP2-SST}

\documentclass{article}
\pdfpagewidth=8.5in
\pdfpageheight=11in

\usepackage{ijcai26}

\usepackage{times}
\usepackage{soul}
\usepackage{url}
\usepackage[hidelinks]{hyperref}
\usepackage[utf8]{inputenc}
\usepackage[small]{caption}
\usepackage{graphicx}
\usepackage{amsmath}
\usepackage{amsthm}
\usepackage{booktabs}
\usepackage{algorithm}
\usepackage{algorithmic}
\usepackage[switch]{lineno}

% Comment out this line in the camera-ready submission
\linenumbers

\urlstyle{same}

\usepackage{color}
\usepackage{multirow}
\usepackage{amssymb}

% To-be-filled placeholders: explicit context wrapped in red asterisks.
\newcommand{\tbd}[1]{\textcolor{red}{*#1*}}

% Please leave this \pdfinfo block untouched for submission and camera-ready.
% Do not include Title and Author information in the pdfinfo section.
\pdfinfo{
/TemplateVersion (IJCAI.2026.0)
}

\title{ProtBLIP2-SST: Protein BLIP-2 Training with Sequence, Structure, and Text}

% Camera-ready author block (uncomment and fill; keep commented for double-blind review):
% \author{
% \tbd{First Author}
% \affiliations
% \tbd{Affiliation 1, City, Country}
% \emails
% \tbd{first.author@institution.edu}
% \And
% \tbd{Second Author}
% \affiliations
% \tbd{Affiliation 2, City, Country}
% \emails
% \tbd{second.author@institution.edu}
% }

\author{Anonymous Authors}

\begin{document}

\maketitle

\begin{abstract}
Bridging protein \textbf{S}equence, \textbf{S}tructure, and \textbf{T}ext (SST) remains difficult for end-to-end generation.
We present \textbf{ProtBLIP2-SST}, a BLIP-2-style framework that keeps a frozen protein encoder, an 8-query Q-Former, and a Galactica-1.3B decoder with mid-layer LoRA fixed, while swapping only the encoder: ESM-2 (sequence-only), SaProt (interleaved amino-acid and Foldseek 3Di tokens), or ProstT5 (bilingual sequence--structure T5).
Stage~1 aligns protein and text via contrastive, matching, and captioning losses; stage~2 generates Swiss-Prot captions.
On an aligned 9{,}239-sample test split, SaProt-650M reaches 33.07\% exact match and 58.53 BLEU-4, modestly above ESM-2-650M (32.83\% / 58.12) with complementary errors: structure conditioning reduces fold-specific enzyme confusions, whereas sequence-only encoding remains stronger on short family-label captions.
We also report parallel cached PLM embeddings that accelerate training by up to $12\times$ on structure-heavy encoders.
\end{abstract}

\section{Introduction}
\label{sec:intro}

Proteins are simultaneously described by amino-acid sequences, three-dimensional structures, and free-text annotations in Swiss-Prot~\cite{bairoch2000swissprot}.
Large language models (LLMs) interpret biomedical text but cannot directly read proteins~\cite{taylor2022galactica}.
Protein language models (PLMs) such as ESM-2~\cite{lin2023esm2} learn sequence representations; AlphaFold2~\cite{varadi2022alphafolddb} further enables structure-centric modeling via SaProt~\cite{su2023saprot} and ProstT5~\cite{heinzinger2024prostt5}.
ProTrek unifies sequence, structure, and function text for retrieval~\cite{su2025protrek}, while ProtT3~\cite{liu2024prott3} adapts BLIP-2~\cite{li2023blip2} for protein-to-text generation.
Prot2Text-V2~\cite{fei2025prot2text} replaces the Q-Former with hybrid sequence-level contrastive alignment (H-SCALE) and a larger decoder.

Three gaps motivate ProtBLIP2-SST:
\textbf{(i)}~structure integration is fragmented across encoders;
\textbf{(ii)}~aggregate caption metrics hide distinct failure modes (pathway swapping vs.\ field hallucination);
\textbf{(iii)}~frozen PLM forward passes dominate training cost.
We keep the ProtT3/BLIP-2 skeleton fixed and vary only the SST encoder, enabling controlled comparisons under identical Q-Former and Galactica configurations.

\textbf{Contributions.}
(1)~A modular BLIP-2 SST framework with interchangeable ESM-2, SaProt, and ProstT5 encoders.
(2)~Aligned Swiss-Prot captioning on 9{,}239 proteins with 114 directional win cases.
(3)~Field-stratified analysis of structure vs.\ sequence error mechanisms.
(4)~Parallel cached PLM embeddings for reproducible multi-encoder training.
\tbd{Optional fifth contribution after ProstT5 stage-2 benchmark completion.}

\section{Related Work}
\label{sec:related}

\paragraph{Protein language models.}
ESM-2~\cite{lin2023esm2} scales masked language modeling on residue sequences.
SaProt interleaves residue and 3Di tokens from Foldseek~\cite{vankempen2024foldseek}.
ProstT5 translates between amino-acid and 3Di sequences with task prefixes~\cite{heinzinger2024prostt5}.

\paragraph{Protein--text modeling.}
Contrastive models (ProtST~\cite{xu2023protst}, ProteinChat~\cite{guo2023proteinchat}) focus on retrieval or linear projection.
ProtT3~\cite{liu2024prott3} introduces two-stage BLIP-2 training for captioning, QA, and retrieval.
Prot2Text-V2~\cite{fei2025prot2text} uses H-SCALE and LLaMA-3.1-8B-Instruct.
ProTrek~\cite{su2025protrek} performs trimodal contrastive alignment for search rather than open-ended generation.

\paragraph{BLIP-2.}
BLIP-2 bridges frozen unimodal encoders and LLMs with a lightweight Q-Former trained for representation learning then generative conditioning~\cite{li2023blip2}.
ProtBLIP2-SST transposes the image encoder to a frozen PLM and the LLM to Galactica (OPT architecture)~\cite{taylor2022galactica,hu2022lora}.

\section{Method}
\label{sec:method}

\subsection{Overview}

ProtBLIP2-SST follows ProtT3's two-stage curriculum (Figure~\ref{fig:architecture}).
A frozen PLM yields per-residue states $\mathbf{H}_p \in \mathbb{R}^{L \times d_p}$.
An 8-query Q-Former cross-attends to $\mathbf{H}_p$, producing $\mathbf{Z}_p \in \mathbb{R}^{8 \times d_q}$.
In stage~2, \texttt{opt\_proj} maps $\mathbf{Z}_p$ into Galactica-1.3B embeddings prepended to prompt and caption tokens; loss applies only to caption spans.

\begin{figure}[t]
  \centering
  % \includegraphics[width=\linewidth]{\tbd{figures/protblip2_sst_architecture.pdf}}
  \fbox{\parbox{0.95\linewidth}{\centering\small
    \tbd{Replace with architecture figure (Stage-1 PTC/PTM/PCap; Stage-2 Galactica+mid-LoRA)}\\
    Frozen PLM $\rightarrow$ Q-Former (8 queries) $\rightarrow$ opt\_proj $\rightarrow$ Galactica-1.3B
  }}
  \caption{ProtBLIP2-SST architecture. Only the PLM backbone changes across SST variants.}
  \label{fig:architecture}
\end{figure}

\subsection{Protein Encoders}

\paragraph{ESM-2 (sequence-only).}
Amino-acid tokenization with \texttt{esm2\_t12\_35M} and \texttt{esm2\_t33\_650M}; PubMedBERT text tower.

\paragraph{SaProt (sequence + structure).}
Interleaved amino-acid and lowercase 3Di tokens from AlphaFold2 via Foldseek~\cite{su2023saprot}; BiomedBERT text tower; checkpoints \texttt{SaProt\_35M\_AF2} and \texttt{SaProt\_650M\_AF2}.

\paragraph{ProstT5 (bilingual SST).}
T5 encoder with \texttt{<AA2fold>} / \texttt{<fold2AA>} prefixes~\cite{heinzinger2024prostt5}.
Stage-1/2 code paths are implemented; full caption benchmark is \tbd{pending ProstT5 stage-2 evaluation on aligned 9239 test split}.

\subsection{Stage 1: Protein--Text Alignment}

Three objectives: \textbf{PTC} (protein--text contrastive), \textbf{PTM} (matching), and \textbf{PCap} (caption LM).
Retrieval uses PTC coarse search with PTM re-ranking over top-128 candidates~\cite{liu2024prott3}.
Stage~1 retrieval numbers on Swiss-Prot are \tbd{to be reported after re-running stage-1 eval for all three encoders}.

\subsection{Stage 2: Caption Generation}

Galactica-1.3B is loaded via \texttt{OPTForCausalLM}.
\texttt{mid\_lora} ($r{=}8$, $\alpha{=}16$, dropout $0.1$) attaches to attention and FFN projections; base weights stay frozen.
Training: 10 epochs, AdamW lr $10^{-4}$, batch size 32, bf16 mixed precision, beam search (beam 3) at inference.

\subsection{Cached PLM Embeddings}

When the PLM is frozen, we precompute hidden states before Q-Former LayerNorm, with multi-GPU sharded cache builders.
LayerNorm is applied online.
Observed speedups on A100: $\sim 2\times$ (ESM-35M), $\sim 6$--$7\times$ (ESM-650M), $\sim 12\times$ (ProstT5 stage-1).
One-time cache build cost: \tbd{report GPU-hours per encoder after final benchmarking}.

\section{Experiments}
\label{sec:experiments}

\subsection{Setup}

\textbf{Data.}
Swiss-Prot V3 from ProtT3~\cite{liu2024prott3}; SaProt uses \texttt{*\_saprot.json}, ESM-2/ProstT5 use \texttt{*\_afdbfilt.json}.
Stage-2 prompt: \texttt{Describe the function of this protein.}
Fair encoder comparison uses an \textbf{aligned 9{,}239-sample} test subset with identical captions and order.
ProtT3's original split has 10{,}000 test proteins~\cite{liu2024prott3}; Prot2Text-V2 reuses that split~\cite{fei2025prot2text}.

\textbf{Metrics.}
Exact match, BLEU-4, ROUGE-L, METEOR~\cite{papineni2002bleu,lin2004rouge,banerjee2005meteor}.
\textbf{Code.}
\tbd{public repository URL to be released upon acceptance}.

\subsection{Main Captioning Results}

Table~\ref{tab:main-caption} reports re-scored stage-2 predictions.
SaProt-650M attains the best metrics among trained variants; gaps to ESM-2 are modest but error types differ (Section~\ref{sec:analysis}).

\begin{table}[t]
  \caption{Swiss-Prot captioning on aligned 9{,}239 test samples.}
  \label{tab:main-caption}
  \centering
  \small
  \begin{tabular}{lccccc}
    \toprule
    Variant & PLM & Ex. & B-4 & R-L & Mtr \\
    \midrule
    SST-ESM   & 35M  & 30.33 & 54.74 & 64.29 & 65.99 \\
    SST-ESM   & 650M & 32.83 & 58.12 & 67.77 & 69.10 \\
    SST-SaProt& 35M  & 30.59 & 55.37 & 65.26 & 66.62 \\
    SST-SaProt& 650M & \textbf{33.07} & \textbf{58.53} & \textbf{68.23} & \textbf{69.54} \\
    SST-ProstT5 & \tbd{T5-encoder size} & \tbd{--} & \tbd{--} & \tbd{--} & \tbd{--} \\
    \bottomrule
  \end{tabular}
\end{table}

\noindent Literature baselines on \emph{different} splits/backbones: ProtT3 (ESM-150M) 25.74\% / 51.47 B-4~\cite{liu2024prott3}; Prot2Text-V2 (ESM-3B + LLaMA-8B) 49.04\% / 57.29 B-4~\cite{fei2025prot2text}.

\subsection{Field-Stratified Analysis}
\label{sec:analysis}

Table~\ref{tab:field-exact} counts exclusive 650M exact-match wins by caption lead field.
SaProt leads on \texttt{FUNCTION:} and \texttt{SIMILARITY:}; ESM-2 leads on \texttt{SUBCELLULAR:}.

\begin{table}[t]
  \caption{Exclusive exact-match wins (650M, $n{=}9{,}239$).}
  \label{tab:field-exact}
  \centering
  \small
  \begin{tabular}{lrrrrr}
    \toprule
    Lead & $n$ & Sa & ESM & Both & None \\
    \midrule
    \texttt{FUNCTION:}    & 8001 & 211 & 175 & 2124 & 5491 \\
    \texttt{SUBCELLULAR:} &  638 &  33 & \textbf{54} &  318 &  233 \\
    \texttt{SIMILARITY:}  &  600 & \textbf{39} &  32 &  330 &  199 \\
    \bottomrule
  \end{tabular}
\end{table}

\textbf{SaProt wins (69 cross-size cases).}
47/69 share \texttt{FUNCTION:} but ESM-2 names the wrong pathway (e.g., lysine biosynthesis vs.\ methanogenesis; tRNA vs.\ mRNA methyltransferase).
SaProt-35M stays exact in 49/69 cases vs.\ 0/69 for ESM-35M, indicating errors in the protein$\rightarrow$Q-Former path.

\textbf{ESM-2 wins (45 cases).}
Short \texttt{SIMILARITY:} family labels from sequence motifs (HSP70, asp23, IS21) favor ESM-2; SaProt sometimes over-generates \texttt{FUNCTION:} prose.
Both models open with \texttt{FUNCTION:} on $\sim$24\% of pure \texttt{SIMILARITY:} targets (shared Galactica prior).

\tbd{Add 2--3 qualitative prediction examples in an appendix figure for enzyme disambiguation and family-label cases.}

\section{Discussion and Conclusion}

ProtBLIP2-SST shows that BLIP-2-style bridging is \emph{encoder-agnostic}: the same Q-Former and Galactica decoder accommodate sequence-only and structure-aware PLMs without redesign.
SaProt improves fold-specific disambiguation; ESM-2 remains competitive on motif-driven family names.
Shared Galactica priors still bias toward verbose \texttt{FUNCTION:} text---future work includes field-balanced sampling and SSF chain-of-thought intermediates~\cite{su2025protrek}.
Cached embeddings make multi-encoder studies practical at scale.

\textbf{Limitations.}
Comparisons mix encoder architecture and input representation (AA+3Di vs.\ AA-only).
ProstT5 caption results and stage-1 retrieval ablations are incomplete.
Decoder scaling (H-SCALE, LLaMA) is not yet explored within ProtBLIP2-SST.
\tbd{Add statistical significance tests and human expert evaluation if required by reviewers.}

We introduced ProtBLIP2-SST, unifying Sequence, Structure, and Text protein captioning through swappable frozen PLMs and a shared BLIP-2 generator, with systematic analysis of when structure conditioning helps and when sequence signal suffices.

\section*{Acknowledgments}
\tbd{Funding sources, compute grants, and collaborator acknowledgments omitted for double-blind review; restore for camera-ready.}

%% Uncomment for camera-ready:
%% \section*{Ethical Statement}
%% \tbd{Ethical statement if required by IJCAI-ECAI 2026.}

\bibliographystyle{named}
\bibliography{main}

\end{document}
